% !TeX root = ../../main.tex

\subsection{Compiler architecture overview}
\label{sec:compiler}

A compiler is a program that translates source code written in one language
into an equivalent form in another language.
When the target and source languages are both high-level languages,
it is also called a transpiler or a source-to-source compiler.
Modern compilation consists of several passes,
each turning one representation into the next.
A compiler then emits target code,
while an interpreter evaluates its internal representation directly,
for example by walking a syntax tree or executing byte code.
Many systems combine both approaches,
for example a runtime may interpret an internal representation
and decide to JIT-compile some portions to native code.\footnote{
	Just-in-time (JIT) compilation translates code
	into native machine code at runtime,
	often to speed up frequently executed parts of a program.
}

\subsubsection{Lexical analysis}

Lexical analysis is normally the first pass of a compiler.
The lexer reads the raw characters and groups them into tokens.
Tokens are indivisible units, each with a kind, its payload
and its position in the source for user-friendly diagnostics.
Whitespace and comments typically carry no meaning and are discarded,
everything else becomes a token:

\begin{figure}[H]
	\centering
	\begin{tikzpicture}[
		tok/.style={draw=black!70, minimum height=0.7cm, inner sep=4pt, outer sep=0, font=\ttfamily\footnotesize, align=center},
		kw/.style={tok, fill=cyan700!12},
		ident/.style={tok, fill=white},
		op/.style={tok, fill=teal700!12},
		num/.style={tok, fill=yellow700!12},
		kind/.style={font=\scriptsize, text=black!65},
	]
	\node[kw] (t1) {SELECT};
	\node[kind, below=1pt of t1] {keyword};
	\node[ident, right=4pt of t1] (t2) {name};
	\node[kind, below=1pt of t2] {identifier};
	\node[kw, right=4pt of t2] (t3) {FROM};
	\node[kind, below=1pt of t3] {keyword};
	\node[ident, right=4pt of t3] (t4) {users};
	\node[kind, below=1pt of t4] {identifier};
	\node[kw, right=4pt of t4] (t5) {WHERE};
	\node[kind, below=1pt of t5] {keyword};
	\node[ident, right=4pt of t5] (t6) {age};
	\node[kind, below=1pt of t6] {identifier};
	\node[op, right=4pt of t6] (t7) {>=};
	\node[kind, below=1pt of t7] {operator};
	\node[num, right=4pt of t7] (t8) {18};
	\node[kind, below=1pt of t8] {number};
\end{tikzpicture}

	\caption{The lexer groups source characters into tokens. Gaps between tokens are discarded whitespace.}
	\label{fig:tokens}
\end{figure}

Keywords such as \mono{SELECT} are recognised by lookup in a fixed table,
while identifiers, numbers and operators are recognised by their shape.
The lexer knows nothing about syntactic structure,
it only rejects input that forms no valid token at all,
such as an unterminated string literal.

\subsubsection{Syntax analysis}

The parser arranges the stream of tokens produced by lexical analysis
into a tree that mirrors the grammatical structure of the input.
A grammar is a set of language rules describing how smaller parts form larger ones,
for example that a comparison consists of two expressions joined by an operator.\footnote{
	Grammars are often written in a standardized notation,
	such as Extended Backus–Naur Form (EBNF),
	which defines a formal syntax to specify a language's syntax rules.
}
The resulting syntax tree drops punctuation that carries no meaning by itself,
such as parentheses and commas, and keeps only the nesting:

\begin{figure}[H]
	\centering
	\begin{tikzpicture}[
		opnode/.style={draw=black!70, fill=teal700!12, rounded corners=3pt, minimum width=1.2cm, minimum height=0.7cm, inner sep=4pt, font=\ttfamily\small, align=center},
		leaf/.style={draw=black!70, fill=white, rounded corners=3pt, minimum width=1.2cm, minimum height=0.7cm, inner sep=4pt, font=\ttfamily\small, align=center},
		level 1/.style={sibling distance=4cm, level distance=1.4cm},
		level 2/.style={sibling distance=2cm, level distance=1.4cm},
	]
	\node[opnode] {+}
	child { node[leaf] {total} }
	child { node[opnode] {*}
			child { node[leaf] {1} }
			child { node[leaf] {2} }
		};
\end{tikzpicture}

	\caption{Syntax tree for \mono{total + 1 * 2}. The shape groups \mono{1 * 2} first, encoding that multiplication binds tighter than addition.}
	\label{fig:ast}
\end{figure}

The tree shape is what concretizes precedence and associativity.
For example, if the grammar defines that multiplication binds tighter than addition,
the parser must nest it deeper.
Or if a repeated operator is defined to group in a fixed direction,
the parser encodes this as nested structure,
e.g.\ \mono{a - b - c} as \mono{(a - b) - c} (left-associative).
A token that fits no rule at its position is a syntax error.
Instead of stopping at the first invalid token,
parsers typically skip ahead to a synchronization point,
such as the end of the line, to collect further errors in the same run.

\subsubsection{Semantic analysis}

A syntax tree can be grammatically valid yet semantically undefined,
it may name things that do not exist or combine values in ways that make no sense.
Semantic analysis checks meaning against the source language's rules.
Name resolution connects each identifier to the declaration it refers to,
following scoping rules: names belong to scopes, scopes nest,
and an inner scope may shadow an outer one.\footnote{
	Shadowing means an inner declaration hides
	an outer declaration of the same name.
	Within the inner scope, the name then refers to the inner declaration.
}
Type checking then verifies that every source construct is applied to
values of suitable types, according to specific rules for each construct.
For example a comparison always yields a boolean,
and arithmetic cannot be performed between a string and a number.

Languages differ in when this checking happens.
A statically typed language checks types before execution
and rejects mismatched programs outright,
whereas a dynamically typed language attaches types to values at runtime
and reports mismatches only when the faulty operation runs.
Static type checking moves a whole class of errors from execution to compilation,
at the cost of requiring enough type information up front.
Whether checking happens statically is independent of how results are produced,
a system can check types up front and still evaluate the tree directly.

\subsubsection{Code generation and interpretation}

Once the tree is checked, there are two main approaches to obtaining results.
Code generation walks the tree and emits target code,
typically one target construct per tree node.
Each node type has a fixed translation pattern,
and the generator applies it recursively from the leaves up.
The mapping is direct when the source and target languages share a construct.
Abstractions for which the target language has no equivalent
may also expand into several target statement.
Interpretation instead evaluates the tree directly.
An executor walks the same nodes and carries out what each one means,
computing values bottom-up without ever producing a standalone target program.
If the tree was type checked beforehand,
both approaches can assume well-typed input and focus purely on producing results,
otherwise runtime type guards are additionally required.

\subsubsection{Errors and diagnostics}

Every compilation pass reports problems it encounters as diagnostics.
A diagnostic typically consists of a message, the source location it refers to,
and a machine-readable error code.
To keep reporting after the first encountered issue, each pass recovers in its own way.
The lexer skips the offending character, the parser jumps to a synchronization point,
and the semantic analyzer may assign a placeholder type
so that checking of independent portions can still be carried out.
Diagnostics distinguish errors, which abort compilation,
from warnings, which flag suspicious but translatable input.
